\documentclass[12pt]{article}
\usepackage[utf8]{inputenc}
\usepackage{graphicx}
\usepackage{amsmath}
\usepackage{amssymb}
\usepackage{mathrsfs}
\usepackage{pgfornament}
\usepackage[many]{tcolorbox}
\usepackage[margin = 0.5in]{geometry}
\usepackage{collectbox}
\usepackage{mdframed}
\usepackage{xcolor}
\usepackage{mathtools}
\DeclarePairedDelimiter\ceil{\lceil}{\rceil}
\DeclarePairedDelimiter\floor{\lfloor}{\rfloor}
\newcommand\norm[1]{\left\lVert#1\right\rVert}
\newcommand\textbox[1]{%
  \parbox{.333\textwidth}{#1}%
}

\linespread{2}

\makeatletter
\newcommand{\mybox}{%
    \collectbox{%
        \setlength{\fboxsep}{2pt}%
        \fbox{\BOXCONTENT}%
    }%
}
\makeatother

\usepackage{listings}
\usepackage{color}

\definecolor{dkgreen}{rgb}{0,0.6,0}
\definecolor{gray}{rgb}{0.5,0.5,0.5}
\definecolor{mauve}{rgb}{0.58,0,0.82}

\lstset{frame=tb,
  language=R,
  aboveskip=3mm,
  belowskip=3mm,
  showstringspaces=false,
  columns=flexible,
  basicstyle={\small\ttfamily},
  numbers=none,
  numberstyle=\tiny\color{gray},
  keywordstyle=\color{blue},
  commentstyle=\color{dkgreen},
  stringstyle=\color{mauve},
  breaklines=true,
  breakatwhitespace=true,
  tabsize=3
}
\title{MAT 523 HW 4}
\begin{document}

\hfill Feb. 27th, 2019
\begin{center}
\begin{large}
\textbf{MAT 524: Functions of a Real Variable II \\
\vspace{-5mm}
Homework IV \\}
\end{large}
\vspace{-2mm}
Nolan R. H. Gagnon \\
\end{center}

\noindent \textbf{{\huge [}{\Large [}} Problem One \textbf{{\Large ]}{\huge ]}}{\Large \textbf{:}}

\vspace{3mm}

\noindent\mybox{\colorbox{yellow}{\textbf{Problem Statement}}}\hspace{3mm}\hrulefill

\noindent Suppose $h\in L^p(E)$ and that $f$ is a measurable function on $E$. If $\norm{h-f}_p<\infty$, show that $f\in L^p(E).$

\vspace{5mm}

\noindent\mybox{\colorbox{yellow}{\textbf{Solution}}}\hspace{3mm}\hrulefill

\noindent Since $h, h-f \in L^p(E)$ and $L^p(E)$ is a vector space,
\begin{align}
    f &= h + (-1)(h-f)
\end{align}
is in $L^p(E)$.

\hfill$\blacksquare$

\vfill

\begin{center}
\pgfornament[width = 50mm, color = black, symmetry = h]{83}\\
\vspace{2mm}
\vspace{-0.65cm}
\hrulefill \\
\vspace{-0.95cm}
\hrulefill
\end{center}

\pagebreak

\noindent \textbf{{\huge [}{\Large [}} Problem Two \textbf{{\Large ]}{\huge ]}}{\Large \textbf{:}}

\vspace{3mm}

\noindent\mybox{\colorbox{yellow}{\textbf{Problem Statement}}}\hspace{3mm}\hrulefill 

\noindent Show that the step functions are not dense in $L^{\infty}([0,1]).$

\vspace{5mm}

\noindent\mybox{\colorbox{yellow}{\textbf{Solution}}}\hspace{3mm}\hrulefill

\noindent Let $A = \bigcup\limits_{n=1}^{\infty}\left[\frac{1}{4n+1},\frac{1}{4n-1}\right]\subset [0,1].$ The characteristic function $\chi_A$ oscillates between $0$ and $1$ infinitely many times, due to the nature of the set $A$. Note that there is no step function that can mimic this behavior, since, by definition, step functions can only change values finitely many times on any interval. With this in mind, let $s$ be any step function on $[0,1]$. On the interval $[0,1/3]$, there must be some interval of positive measure where the distance between $s$ and $\chi_A$ is greater than or equal to $0.5$, i.e., $\norm{\chi_A - s}_{\infty}\geq 0.5$. Thus, any $\varepsilon$ challenge with $\varepsilon \leq 0.5$ will fail, meaning the step functions are not dense in $L^{\infty}([0,1]).$   

\hfill$\blacksquare$

\vfill

\begin{center}
\pgfornament[width = 50mm, color = black, symmetry = h]{83}\\
\vspace{2mm}
\vspace{-0.65cm}
\hrulefill \\
\vspace{-0.95cm}
\hrulefill
\end{center}

\end{document}
